\documentclass[11pt,letterpaper]{article}
\usepackage[margin=0.88in,headheight=15pt]{geometry}
\usepackage[T1]{fontenc}
\usepackage{lmodern,microtype,amsmath,amssymb,amsthm,mathtools}
\usepackage{booktabs,tabularx,array,xcolor,fancyhdr}
\usepackage[colorlinks=true,linkcolor=blue!45!black,urlcolor=blue!50!black,citecolor=blue!45!black]{hyperref}
\definecolor{ink}{RGB}{28,49,68}
\pagestyle{fancy}\fancyhf{}
\setlength{\parindent}{0pt}\setlength{\parskip}{5pt}
\setlength{\emergencystretch}{2em}
\numberwithin{equation}{section}
\theoremstyle{plain}
\newtheorem{theorem}{Theorem}[section]
\newtheorem{lemma}[theorem]{Lemma}
\newtheorem{proposition}[theorem]{Proposition}
\theoremstyle{remark}\newtheorem{remark}[theorem]{Scope}
\DeclareMathOperator{\Cov}{Cov}\DeclareMathOperator{\supp}{supp}
\DeclareMathOperator{\Part}{Part}\DeclareMathOperator{\Ad}{Ad}
\DeclareMathOperator{\dist}{dist}\DeclareMathOperator{\diag}{diag}
\newcommand{\dd}{\,\mathrm d}\newcommand{\eps}{\varepsilon}
\newcommand{\norm}[1]{\left\lVert#1\right\rVert}
\newcommand{\abs}[1]{\left|#1\right|}
\newcommand{\R}{\mathbb R}\newcommand{\C}{\mathbb C}
\title{Boundary traces and reconstruction at a changing base}\author{}\date{}
\begin{document}\maketitle
\section{The actual target and two preliminary observations}

For a fixed original history $\omega$ and the complete simultaneous wall
tuple $\mathbf s$, keep
\begin{equation}
 W_{\omega,\mathbf s}
 =P_{\rm old}\bigl(q_{\rm old}\Gamma^R_{\omega,\mathbf s}\bigr),
 \qquad
 F_{j,e;\omega,\mathbf s}
 =R_j\bigl(W_{\omega,\mathbf s},J_{\rm phys}(W_{\omega,\mathbf s})\bigr).
 \label{eq:stored}
\end{equation}
Every expression in \eqref{eq:stored} is used only where its original map
is defined. The map $q_{\rm old}$ includes the complete original determining
record: oriented averages, bare identity rows, meshes, path prefixes,
endpoint frames, exterior values and boundary conditions. It is not just
the restriction to the source's geometric support. $P_{\rm old}$ is its
original constrained Wilson minimum.

For the clean-born $d=2$ stable $H$ packet, write
\begin{equation}
 r=L^{j-e},\qquad \rho=g_e/g_j,\qquad u=g_e/\alpha,
 \qquad (\gamma_2,\gamma_3,\gamma_4)=(s,b,2),
 \label{eq:notation}
\end{equation}
where $b$ is the cubic age exponent, not the Wilson inverse coupling.
The already proved last-regular envelope has the form
\begin{equation}
 |F|\le A_j\rho^2
 \left[r^sQ^2+ur^bQ^3+u^2r^2Q^4\right],\qquad Q=q_H\ge0.
 \label{eq:envelope}
\end{equation}
Its application to \eqref{eq:stored} is the missing issue. In particular,
$Q$ includes pinned and exterior determining data as well as integrated
fluctuations. Formula \eqref{eq:envelope} cannot be assumed on a new mixed
domain simply because it holds on the last admitted regular problem.
The saved current/source derivatives and anchor norms are additional
finite rows of the packet, not consequences of a scalar value bound.

The first observation holds: squaring
\eqref{eq:envelope} needs moments through degree eight, and the full physical
core can lose canonical coupling factors. Sections 5--8 show which such
losses an existing reserve can pay. The second observation holds at first
order on its expressly admitted protected chart. Section 4 separates that
fact from the stronger higher-jet and anchor requirements.

All results below concern this fixed stored object. A source-only
reconstruction agrees with the physical source at the genuine minimum
endpoint as in LM46, but endpoint equality alone supplies neither
simultaneous-wall admission nor a new source-independent physical radius.

\section{Retained targets and joined boundary traces}

\begin{lemma}[Exact old-target retention]
On the admitted domain of the complete old determining problem,
\begin{equation}
 q_{\rm old}P_{\rm old}(v)=v.
 \label{eq:rightinverse}
\end{equation}
Consequently, for every admitted parameter family $\Gamma(z)$,
\begin{equation}
 q_{\rm old}\bigl(P_{\rm old}q_{\rm old}\Gamma(z)\bigr)
 =q_{\rm old}\Gamma(z).
 \label{eq:targetidentity}
\end{equation}
All finite ordered derivatives of these composite target functions agree.
The statement retains the specified common gauge or endpoint conjugations.
\end{lemma}
\begin{proof}
Equation \eqref{eq:rightinverse} is the defining minimum constraint.
Substitute $v=q_{\rm old}\Gamma(z)$ and differentiate the identity.
For a nonlinear row the first derivative reads
\[
 Dq_{\rm old}(W)DW=Dq_{\rm old}(\Gamma)D\Gamma;
\]
it is not $q_{\rm old}DW=q_{\rm old}D\Gamma$. Higher derivatives likewise
differentiate the entire composite.
\end{proof}

In particular, admission to the saved spare domain requires
\begin{equation}
 q_{\rm old}\Gamma^R_{\omega,\mathbf s}
 \in q_{\rm old}\bigl(\pi_U\mathcal C^{\rm spare}_{j,e}\bigr)
 \quad\hbox{on the full original simultaneous domain}.
 \label{eq:necessary}
\end{equation}
Here $\pi_U$ is the field projection of the saved paired domain.
This is necessary, not sufficient. Reconstruction can restore the
own-current relation once admitted; it cannot repair targets outside this
image.

There is a precise current consequence. On every old pinned identity bond,
$W_b=\Gamma_b$. A Wilson plaquette word is therefore unchanged when its
whole stencil is pinned. The physical Wilson current at a bond is unchanged
only when \emph{all} incident-plaquette bond stencils are pinned, with the
original action coefficients. Every finite jet of those complete words is
also unchanged. Hence every nonnegative weighted aggregate of those target
or pinned-stencil jets is exactly unchanged. An interior averaging gain
cannot supply their missing boundary or anchor bound. This is not an
example of a bad selected history.

\subsection{The base term cannot be discarded}

In one admitted target chart, a join admits the exact bookkeeping identity
\begin{equation}
 q\Gamma_i-q\Gamma_{i-1}
 =\bigl(q\Gamma_i-q\Gamma_{i,0}\bigr)
  +\bigl(q\Gamma_{i,0}-q\Gamma_{i-1}\bigr).
 \label{eq:base}
\end{equation}
The first term responds to the localized correction. The second contains
the joined base minimum and its re-averaged boundary trace. A small estimate
for the first term does not estimate the second. This is an identity in a
fixed chart, not an assumption that group-valued averages are additive.

R-II (1.50) retains the genuine old constrained minimum at the physical
endpoint. Its later joins have distinct coincidence interiors $Y_i^0$ and
changed complementary boundary rows. The gauge in (1.61) is retained.
The range in (1.65) is the original native member range. The
source does not identify it with $\mathcal C^{\rm spare}_{j,e}$.
The discussion agrees with LM42--48 and PE12 of~\cite[pp.~781--783, 1174--1180]{W} and with~\cite{rii}.

The actual localization construction also does not supply the needed
universal collar. $Y_1$ is a smallest containing localization domain, and
later added components connect to successive boundary layers. Containment
of the complete old readset in $Y_i$ does not imply its positive separation
from $Y_i\setminus Y_i^0$. The Green forcing uses complete row-adjoint
stencils, including prefixes, not just row labels. III's join construction
does not state the missing saved-readset erosion~\cite{iii}.

For genuine surviving ancestors there is a restricted positive statement.
If the complete packet readset is unchanged, its first touch is $e$, and
the verified survival step is $Z_{e-1}=S(Z_{e-2})$, the original ten-layer
padding gives
\begin{equation}
 \dist_\infty(S_{\rm old},Z_{e-2})
 \ge cM R_{e-1}a_{e-1}.
 \label{eq:ancestorcollar}
\end{equation}
Here $a_{e-1}$ is the physical scale length, not an ancestry index.
An adjoint stencil contained in that actual ancestor inherits this
separation. Artificial join rows, moving readsets, and native-metric depth
estimates are not supplied by \eqref{eq:ancestorcollar}.

\section{The exact reconstruction response, including the changed base}

These identities use a common finite-dimensional gauge-fixed chart of an
\emph{already admitted} old minimum. The original nonredundant target and
gauge quotient are retained, and its augmented derivative is assumed
invertible on the verified response domain. Admission or uniqueness alone
would not imply that nondegeneracy. On complex walls, the symbol $*$ below
means the bilinear transpose, continued holomorphically from the real
adjoint; it is not complex conjugation. These formulas are not an
existence argument for the chart. Write its stationary equations as
\begin{equation}
 \nabla S(W)-Dq(W)^*\lambda=0,\qquad q(W)=v.
 \label{eq:kkt}
\end{equation}
Set
\begin{equation}
 Q_W=Dq(W),\qquad
 L_W=S''(W)-\sum_\alpha\lambda_\alpha D^2q_\alpha(W),\qquad
 K_W=\begin{pmatrix}L_W&-Q_W^*\\Q_W&0\end{pmatrix}.
 \label{eq:lagrangian}
\end{equation}
The nonlinear-constraint term in $L_W$ must remain. Let $G_W$ and $P_W$
denote the upper-left and upper-right blocks of $K_W^{-1}$; $P_W$ here is
the derivative lift, distinguished from the nonlinear $P_{\rm old}$.
Block multiplication gives
\begin{equation}
 Q_WG_W=0,\quad G_WQ_W^*=0,\quad
 G_WL_W+P_WQ_W=I,\quad Q_WP_W=I.
 \label{eq:block}
\end{equation}

\begin{proposition}[Off-base derivative identity]
For $W=P_{\rm old}q\Gamma$ in this fixed chart,
\begin{align}
 DW&=P_WQ_\Gamma D\Gamma,\label{eq:exactresponse}\\
 DW-D\Gamma
 &=-G_WL_WD\Gamma
   +P_W(Q_\Gamma-Q_W)D\Gamma.\label{eq:offbase}
\end{align}
\end{proposition}
\begin{proof}
Differentiate \eqref{eq:kkt}. The target derivative is $Q_\Gamma D\Gamma$.
Apply the two upper blocks of the inverse, then use \eqref{eq:block}.
\end{proof}

At a common endpoint $W=\Gamma$, choose any smooth extension of the old
multiplier and put
$d_{\rm old}=\nabla S(\Gamma)-Dq(\Gamma)^*\lambda_{\rm ref}$.
Then \eqref{eq:offbase} reduces to
\begin{equation}
 DW-D\Gamma=-G_WDd_{\rm old}.
 \label{eq:defect}
\end{equation}
The arbitrary multiplier derivative drops out by $G_WQ_W^*=0$.
For an independent preparation or source wall, $Dd_{\rm old}$ need not
vanish. Away from the endpoint the second term of \eqref{eq:offbase}
cannot be omitted.

For example, take $S(x,y)=(x^2+y^2)/2$ and $q(x,y)=x+y^2/2$ near
$v<1$. The local constrained minimum is $W=(v,0)$. At
$\Gamma=(v-h^2/2,h)$ with tangent $(c,d)$,
$DW-D\Gamma=(hd,-d)$. The first term of \eqref{eq:offbase} gives
$(0,-d)$ and its second term gives $(hd,0)$. This finite-dimensional
example checks the missing term; it is not an original Wilson-history
counterexample.

If $S_t,q_t$ themselves depend on a wall, define
\[
 b_W=\partial_t\bigl(\nabla S_t-Dq_t^*\lambda\bigr)\big|_{W,\lambda}.
\]
Differentiation yields the additional explicit terms
\begin{equation}
 \dot W=-G_Wb_W+P_W\bigl[
 Q_\Gamma\dot\Gamma+\partial_tq_t(\Gamma)-\partial_tq_t(W)\bigr].
 \label{eq:moving}
\end{equation}
Physical recovery by a moving chart $\chi_t$ also contributes
$D\chi_t(W)\dot W+\partial_t\chi_t(W)$. Boundary and frame responses
cannot be hidden in a fixed-chart identity.

\subsection{Only genuinely shared rows cancel}

At a common flat quadratic center, let $H$ be the positive gauge-fixed
action Hessian and $Q_o,Q_n$ the complete old and new row matrices on their
nonredundant full-row-rank target quotients. Define
\begin{align}
 \mathcal P_o&=H^{-1}Q_o^*(Q_oH^{-1}Q_o^*)^{-1},\\
 \mathcal G_o&=H^{-1}-\mathcal P_oQ_oH^{-1}.
\end{align}
If $\Gamma=\mathcal P_n v$ and $H\Gamma=Q_n^*\mu$, then
\begin{equation}
 (\mathcal P_oQ_o-I)\Gamma=-\mathcal G_oQ_n^*\mu.
\end{equation}
For a verified factorization $Q_n=TQ_o+E$, this becomes
\begin{equation}
 (\mathcal P_oQ_o-I)\Gamma=-\mathcal G_oE^*\mu.
 \label{eq:rowdefect}
\end{equation}
All unmatched boundary, prefix, averaging and transport adjoints remain
in $E^*\mu$. The three-join construction of~\cite{iii} does not supply the
complete $E$, its weighted multiplier sum, or separation of all its
stencils from the saved readset. Energy contraction of $\mathcal P_oQ_o$
does not give the required local $r^2$ field or $r^3$ current response.

Finally, constrained stationarity is not $J=0$. In coordinates,
$\nabla S(W)=Q_W^*\lambda$, and its derivative contains both
$Q_W^*D\lambda$ and $\sum\lambda_\alpha D^2q_\alpha(W)DW$.
Converting this coordinate gradient to the physical Wilson current also
differentiates the actual chart cotangent map and retains the original
gauge-quotient identification. Those boundary reactions
are part of the current estimate still required.

\section{Analysis of protected rescaling, higher jets and anchors}

On the \emph{same admitted protected current chart}, assume the genuine
response bound
\begin{equation}
 \norm{DA[v]}_{C^{1,\beta}}+\norm{DJ_{\rm phys}[v]}_\infty
 \le Cg_e\norm v,
 \label{eq:protectedresponse}
\end{equation}
and the original small background bounds
$\norm A_{C^{1,\beta}}+\norm J_\infty\le\eps$.
The entire enlarged old carrier, its difference and current stencils,
and the neighborhoods used in the Holder seminorm must lie in this same
chart. Take a fixed carrier constant $D$ and $a\le r\le1$, with old mesh
$\xi=a/r$. The current response is a
separate input; it does not follow from $C^1$ control of $A$ alone.
For an old contractible carrier of relative size $r$, use the DA2 gauge
\[
 h(x)=\exp\left(\sum_i(x_i-x_{0,i})A_i(x_0)\right),\qquad
 A_{\rm old}(y)=rA^h(x_0+ry).
\]
The cancellation of the constant connection in the link word gives
\begin{equation}
 \norm{DA_{\rm old}[v]}_{C^{1,\beta}}
 \le C_Dr^2g_e\norm v.
 \label{eq:oldfield}
\end{equation}
Indeed $DA^h$ has the leading difference
$DA(x)-DA(x_0)$ and analytic remainders with the same factor $r$ in
supremum norm; rescaling contributes the second factor. Exact link-word
divisibility is retained before division by the fine mesh.

Using the same fixed Wilson-current and action normalization, without a
running-coupling conversion, the current identity and its derivative are
\begin{align}
 J_{\rm old}&=r^3\Ad_hJ,\\
 DJ_{\rm old}[v]
 &=r^3\Ad_h\bigl(DJ[v]+[h^{-1}Dh[v],J]\bigr),
 \qquad \norm{h^{-1}Dh[v]}\le C_Drg_e\norm v.
 \label{eq:oldcurrent}
\end{align}
Thus $\norm{DJ_{\rm old}[v]}_\infty\le C_Dr^3g_e\norm v$ on that chart.
Complete scalar Ward invariance removes a simultaneous vertical tangent;
it does not delete the displayed commutator while boundary, current or
reference slots are held fixed.

Response estimates leave the base $A(0),J(0)$ intact. The joined bases
are generally not flat when only an internal fluctuation is set to zero.
Those terms must remain in $Q$ or receive a separate bound.

\subsection{Higher coupling grades need the original physical radius}

A first canonical derivative bound is insufficient. On a fixed canonical
complex ball,
\begin{equation}
 A_g(z)=g(z+z^2),\qquad F(A)=g_j^{-2}A^2
 \quad\Longrightarrow\quad
 DA_g=O(g),\qquad (F\circ A_g)'''(0)=12g^2g_j^{-2}.
 \label{eq:jetcounter}
\end{equation}
There is no uniform cubic $O(g^3g_j^{-2})$ conclusion. This is a countermodel
to an inference, not a claim about an original selected history.

A sufficient stronger hypothesis is holomorphic dependence on the
physical variable $w$, with a uniform bound on the first derivative
$D_wP$ on the larger \emph{unchanged} physical radius $\delta$.
Cauchy on a fixed smaller part of that domain gives
\begin{equation}
 \norm{D_z^mP(g_ez)}\le C_mg_e^m\delta^{1-m},\qquad 1\le m\le m_*.
 \label{eq:cauchy}
\end{equation}
The factors involving a shrinking old radius are retained. A first bound
on a fixed canonical ball only gives $O(g_e)$ for higher derivatives.
The full chain rule contains every block
\begin{equation}
 D_I(F\circ\Psi)
 =\sum_{\pi\in\Part(I)}D^{|\pi|}F(\Psi)
      [D_B\Psi:B\in\pi].
\end{equation}
The actual-current hierarchy also contains
$d_a^*d_aD_IA$ in addition to the derivatives of its analytic nonlinear
part. The same simultaneous walls must be admitted throughout the Cauchy
discs; native analyticity on a different determining problem is inadequate.

If an amplitude loses $r^{-\eta}$, the value-envelope powers become
\begin{equation}
 s-2\eta,\qquad b-3\eta,\qquad 2-4\eta.
\end{equation}
The proposed condition
$\eta<\min(s/2,b/3,1/2)$ gives positive remaining exponents. It does not
preserve the exact saved age, prove higher-jet grades, or check old-radius
admission. No such expenditure is silently made here.

\subsection{Absolute anchor sums remain a separate estimate}

A scalar first-order $\ell^\infty$ operator norm equals its unweighted
coefficient $\ell^1$ norm. Weighted locality and higher-order absolute
sums require more. For a symmetric Hadamard matrix $H_n$ with
$H_nH_n^*=nI$, the bilinear form
\[
 B_n(v,w)=n^{-3/2}v^TH_nw
\]
has operator norm at most one on two unit $\ell^\infty$ balls, while the
absolute coefficient sum is $\sqrt n$. Thus ordinary multilinear bounds
and a common Cauchy domain do not supply a uniform absolute anchor norm.

Localized response columns can supply such sums on their admitted domains.
For example, on a uniform mesh a column bound
$|\partial_a\Psi(x)|\le Ce^{-\kappa d(x,a)}$, uniform on a larger physical
disc, gives finite-order kernels bounded by
$C_m\delta^{1-m}e^{-\kappa\max_i d(x,a_i)}$ by Cauchy and symmetry.
Uniform polynomial ball counts then sum any fixed polynomial anchor moment.
For the actual multiscale record, the input measures, scale powers, row
multiplicities and native metric must all be checked. The verified native
response theorem is not a statement for the changed joined target tuple.

\section{The exact moment requirement and the pinned marginal}

\begin{proposition}[Raw three-row moment bound]
On any finite positive body $\mu$ where \eqref{eq:envelope} is actually valid,
\begin{equation}
 \norm F_{L^2(\mu)}\le A_j\rho^2\left[
 r^s(\mu Q^4)^{1/2}+ur^b(\mu Q^6)^{1/2}
 +u^2r^2(\mu Q^8)^{1/2}\right].
 \label{eq:l2}
\end{equation}
\end{proposition}
\begin{proof}
Apply Minkowski to the three nonnegative majorant terms. Squaring their
sum directly gives powers $Q^4,Q^5,Q^6,Q^7,Q^8$; the quartic square
requires degree eight. No artificial factor three is needed in
\eqref{eq:l2}.
\end{proof}

If $Q^2\le c_H(1+E)$ for a genuinely owned energy and
$\mu e^{\theta E}\le B$, define
\begin{equation}
 M_m(\theta)=\sup_{t\ge0}(1+t)^me^{-\theta t}
 =\begin{cases}
 1,&m\le\theta,\\
 (m/\theta)^me^{-m+\theta},&m>\theta.
 \end{cases}
 \label{eq:mm}
\end{equation}
The derivative locates the maximum at $0$ or $m/\theta-1$. Consequently
each $(\mu Q^{2m})^{1/2}$ is at most
$c_H^{m/2}M_m(\theta)^{1/2}\sqrt B$. This uses \emph{one} reserve;
the degree-eight field moment is a degree-four energy polynomial.

\subsection{An internal reserve does not remove fixed determining data}

Disintegrate only when the original body actually does:
$\dd\mu=\dd\nu(\lambda)\dd\mu_\lambda(\xi)$,
$M_\lambda=\mu_\lambda1$. Suppose
$Q\le h(\lambda)+\sqrt{cE_\lambda}$ and
$\mu_\lambda e^{\theta E_\lambda}\le B_\lambda$.
Binomial expansion and Minkowski give, for $m=2,3,4$,
\begin{align}
 \norm{Q^m}_{L^2(\mu_\lambda)}
 \le {}&h^m\sqrt{M_\lambda} \nonumber\\
 &+\sqrt{B_\lambda}\sum_{\ell=1}^m
 \binom m\ell h^{m-\ell}c^{\ell/2}
 \left(\frac{\ell}{\mathrm e\theta}\right)^{\ell/2}.
 \label{eq:pinned}
\end{align}
The leading pinned term keeps the actual mass. If the original larger body
integrates $\lambda$, the quartic row needs, among its terms,
\begin{equation}
 \int h(\lambda)^8M_\lambda\dd\nu(\lambda)\le CB_X.
 \label{eq:pinnedtarget}
\end{equation}
If $\lambda$ remains external, a uniform bound is required instead. There
is no exterior integral or reserve to invent.

For $M=\mu1>0$, put $f_\lambda=\mu_\lambda F/M_\lambda$ and similarly
$g_\lambda$, and $\bar f=\mu F/M$, $\bar g=\mu G/M$. Zero-mass fibers
contribute zero. The full raw covariance is exactly
\begin{align}
 \mu(FG)-\frac{\mu F\,\mu G}{M}
 ={}&\int\left[\mu_\lambda(FG)
       -\frac{\mu_\lambda F\,\mu_\lambda G}{M_\lambda}\right]\dd\nu
       \nonumber\\
 &+\int M_\lambda(f_\lambda-\bar f)(g_\lambda-\bar g)\dd\nu.
 \label{eq:totalcov}
\end{align}
No conditional Poincare or contact prescription deletes the second term.

\subsection{The uniform filter alone has the critical four-dimensional norm}

For the complete uniform $Q_2$ row established in~\cite[\S2]{UA},
\[
 \sum_yK_r(x,y)=1,\quad 0\le K_r(x,y)\le C(a/r)^4,
 \quad |\supp K_r(x,\cdot)|_a\le Cr^4.
\]
Its exact physical $L^2_a$-to-one-value norm is
\begin{equation}
 \left(\sum_yK_r(x,y)^2/a^4\right)^{1/2},
 \qquad cr^{-2}\le\norm{K_r(x,\cdot)}_{(L^2_a)^*}\le Cr^{-2}.
 \label{eq:critical}
\end{equation}
The upper bound uses the maximal coefficient, and the lower bound uses
normalization and support volume. Thus a global curvature-energy estimate
can give $r^2|Q_2F|\le Cg_e\sqrt E$, but it does not supply the extra $r^2$
needed by the protected canonical readout. A sufficient fourth local-energy
moment for the corresponding degree-eight target would be
\begin{equation}
 \int\left(\sum_{\supp K_r}a^4|F|^2\right)^4\dd\mu
 \le Cg_e^8r^{16}B_X.
 \label{eq:localenergytarget}
\end{equation}
This is a route-specific target, not a necessity theorem excluding signed
cancellation. No selected original history is asserted to saturate
\eqref{eq:critical}. Identity and mixed boundary rows do not inherit its
uniform averaging hypothesis.

\section{A summed estimate for genuinely owned native-depth records}

This section improves the numerical part of the aggregate problem without
assuming a separate charged array at every scale. It uses the original
source-free positive recurrence within its existing scope.

Let $\mathfrak V_\omega$ be the set of actual pairs $(t,Z)$ consisting of
an RG step and a live component in the unfolded ancestry contributing to
the completed component $X$. Define
\begin{equation}
 T_\omega=|\mathfrak V_\omega|,\qquad
 \mathcal G_\omega=\sum_{(t,Z)\in\mathfrak V_\omega}
      M^4R_t^5\bigl(d'_t(Z)+1\bigr).
 \label{eq:livet}
\end{equation}
With the fixed original $M,R_t\ge1$, $T_\omega\le\mathcal G_\omega$.
The summands are the original per-component cost majorants LM21.
The measure below is the complete original positive reference body, or
its already justified positive dominating body, with its inherited
scalars, independent normalizers and existing score allocations retained.
It is not a newly normalized conditional fiber.

\begin{proposition}[Fixed live-step exponential weight]
Within the original action, record and history assignments, the positive
completed-body estimate can retain, for any fixed $\lambda\ge0$,
\begin{equation}
 \sum_\omega\int
 e^{\theta E_\omega+\tau N_\omega^{\log}+\lambda T_\omega}
 \dd\mu_\omega
 \le B_{X,\lambda},\qquad
 B_{X,\lambda}\le C_\lambda
 e^{-c_\lambda P_0(x_e)-\kappa_\lambda d(X)},
 \label{eq:timebank}
\end{equation}
for some fixed remaining $\theta,\tau,c_\lambda,\kappa_\lambda>0$ after
all prescribed finite costs are combined once and the weak floor is chosen
accordingly. Here $P_0(x)=(\log x)^{p_0}$ and $x_e=g_e^{-2}$.
\end{proposition}
\begin{proof}
Multiplication by $e^{\lambda T_\omega}$ is bounded by increasing the
coefficient $C_w$ of the existing charge
$C_wM^4R_t^5(d'_t(Z)+1)$ by $\lambda$. Its future sum remains bounded by
$C(C_w+\lambda)M^4L^5R_t^6(d'_t(Z)+1)$.
The original inequalities $p_0>6r_0$ and
$2p_1-9r_0>p_0$ make birth, renewal and surplus merger credits dominate
this enlarged fixed coefficient beyond one weak floor. A no-new-failure
step subtracts its first future charge; a renewal funds the new future
bank; a merger keeps one reserve and pays the overlap from surplus child
reserves. Original history entropy is already in the charge.
This is the same recurrence as (8.3)--(8.4), LM21--23 and S86 of~\cite[pp.~137--138, 774, 614]{W},
with a fixed additional numerical cost. It neither changes the
integral nor assigns energy to new coordinates.
\end{proof}

For a retained record $R$, let $\ell_R$ be its native scale and $a_R$ the
start of its continuous live retention in this same unhealed ancestry.
If the literal chain exists at every step $a_R,\ldots,e-1$, then
$e-a_R\le T_\omega$. Only if $a_R=\ell_R$ does this directly imply
\begin{equation}
 \sigma_R^{-D}=L^{D(e-\ell_R)}
 \le e^{D(\log L)T_\omega},\qquad \sigma_R=L^{\ell_R-e}.
 \label{eq:depth}
\end{equation}
Otherwise the unpaid factor $L^{D(a_R-\ell_R)}$ remains. This is not a
mark's pre-entry regular age to be charged retrospectively.

\begin{proposition}[Whole owned-row aggregate]
Fix $D,m\ge0$ and $\Lambda_e>0$. Let the actual nonnegative weights
of a verified channel obey
$\sum_Rw_{xR}\le C_w'$, and assign its inputs once to the eligible chains
of \eqref{eq:depth}. Suppose the actual native quantity being filtered has
a proved bound
$|c_R|\le C\Lambda_e(1+E_R)^{m/2}$, with $\sum_RE_R\le E_\omega$.
Then
\begin{equation}
 \mathcal U_D(x)=\sum_Rw_{xR}\sigma_R^{-D}|c_R|
 \le CC_w'\Lambda_e e^{D(\log L)T_\omega}(1+E_\omega)^{m/2}.
 \label{eq:aggregatesup}
\end{equation}
For fixed $q\ge1$, choosing $\lambda>qD\log L$ in the same preallocated
recurrence gives
\begin{equation}
 \sum_\omega\int
 \left(\frac{\sup_x\mathcal U_D(x)}{\Lambda_e}\right)^q
 \dd\mu_\omega\le C_{q,D,m}B_{X,\lambda}.
 \label{eq:aggregatemoment}
\end{equation}
\end{proposition}
\begin{proof}
Bound the maximum eligible native depth by \eqref{eq:depth}, then sum the
whole row using $\sum_Rw_{xR}\le C_w'$. The remaining polynomial in
$E_\omega$ is paid by its fixed positive fraction in
\eqref{eq:timebank}. Actual multiplicity polynomials require a proved
bound by the original count ledger; no artificial collar count is assumed.
The fixed constants also depend on the native-envelope constant, $C_w'$
and the available energy fraction.
\end{proof}

For curvature $D=2$. For the directional-current channel $D=3$,
\emph{provided its actual directional derivative envelope is known}.
One waiting coefficient exceeding $24\log L$ therefore handles their
joint eighth moment in this eligible sector. An $O(1)$ native amplitude
normalized by $g_e$ additionally costs $x_e^4$, payable from a remaining
birth rare coefficient as in Section 7.

The qualification about current is essential. A bound on the physical
divergence $J=D^*F$ alone does not bound each $D_\rho F$ used by the
$Q_3$ filters. The elementary bound
$|D_{\rho,a}F|\le2a^{-1}\sup|F|$ retains the actual mesh $a$; it cannot
be replaced by $2\sigma^{-1}\sup|F|$ on a reconstructed patch with
$a\ll\sigma$. Same-mesh compact words with a complete shifted stencil
are a valid restricted case. The general changed-record directional and
current-jet bounds remain open.

Eligible records include genuinely retained own-spacing variables with
their continuous original ancestry. A healed variable is an inherited
coefficient, not a compact variable to reopen. Exterior clean inputs,
preparation subdepth inside a single RG step, and artificial joined
boundary rows have not been assigned to these chains. Thus
\eqref{eq:aggregatemoment} controls an owned subset jointly; it does not
complete \eqref{eq:pinnedtarget} for the full old boundary.

\section{Paying cutoff and full-core losses with one original rare factor}

The numerical argument in (8.3)--(8.4), S86 and LM21--23 of~\cite[pp.~137--138, 614, 774]{W} retains a
positive rare coefficient after the original complete history sum and
fixed identified cost increases. It keeps the birth reserve separate from
the future waiting bank. R-II's original rare-factor recurrence is the
source-free input~\cite[pp.~137--138]{W},~\cite{rii}. Use the notation
\begin{equation}
 B_X^{(c)}=C_0e^{-cP_0(x_e)-\kappa d(X)},\qquad
 P_0(x)=(\log x)^{p_0},\qquad p_0>1.
 \label{eq:bc}
\end{equation}
Its coefficient must actually survive the combined allocation. An
unspecified symbol $B_X$, or a small value of $T_X1$ alone, is insufficient.

For every fixed $m>0$, $\delta>0$ and $x\ge1$, the explicit scalar bound
\begin{equation}
 x^m e^{-\delta(\log x)^p}
 \le \exp\left[
 \frac{p-1}{p}m\left(\frac{m}{\delta p}\right)^{1/(p-1)}
 \right],\qquad p>1,
 \label{eq:rarepower}
\end{equation}
follows by maximizing $mt-\delta t^p$ over $t=\log x\ge0$.
Consequently one fixed reduction $c\mapsto c-\delta$ pays any specified
finite product of current inverse-coupling and logarithmic cutoff powers.
All powers for a finite packet are combined before choosing this reduction.

This is stronger than the minimal VC logarithmic record comparison.
That comparison has $w_e=1+\log(1+x_e)\le C\mathcal A_\omega$ when a
literal earlier charged record and the actual running tube provide it;
it pays fixed polynomials in $w_e$ from a small spare
$e^{\theta\mathcal A_\omega}$. It does not automatically pay arbitrary
$x_e^m$. The remaining birth coefficient in \eqref{eq:bc} is what pays
those inverse-coupling powers.

If only $Q\le CP_e$ is proved on the whole determining record,
\eqref{eq:envelope} leaves
\[
 CA_j\rho^2P_e^2
 \left[r^s+ur^bP_e+u^2r^2P_e^2\right].
\]
The cutoff cost is real until paid. If $Q=u^{-1}=\alpha/g_e$, the three
rows are exactly
\begin{equation}
 A_j(\alpha/g_j)^2r^s,\qquad
 A_j(\alpha/g_j)^2r^b,\qquad
 A_j(\alpha/g_j)^2r^2.
 \label{eq:fullcore}
\end{equation}
These are upper-envelope losses; no saturation or divergence is claimed.

\begin{proposition}[Full-core raw $L^2$ carrying after actual admission]
Suppose the actual full selected support has
\eqref{eq:envelope}, $Q\le C/u$, and its source-free positive mass is
bounded by an identified $B_X^{(c)}$. Keep the given clean radius $\alpha$
bounded above by its original fixed constant. Then for $0<\delta<c$,
\begin{equation}
 \norm F_{L^2(\mu)}\le C A_j\rho^2
 \left[r^s+ur^b+u^2r^2\right]\sqrt{B_X^{(c-\delta)}}.
 \label{eq:fullcorel2}
\end{equation}
\end{proposition}
\begin{proof}
The three terms of \eqref{eq:l2} require the scalar payments
$u^{-4}B^{(c)}$, $u^{-6}B^{(c)}$ and $u^{-8}B^{(c)}$.
The largest fixed inverse-coupling power is
$u^{-8}=\alpha^8x_e^4$. Apply \eqref{eq:rarepower} once to the finite
maximum and retain each original row coefficient outside the payment.
If a factor is less than one it needs no payment.
\end{proof}

This conclusion recovers the displayed $u$ and $u^2$, not just the three
age exponents. It requires the full-core \emph{saved} envelope, not an
unaged native supremum. A rare bound does not define a source outside its
analytic domain, prove its canonical jets, or manufacture an arbitrary
earlier age. R-origin radius costs and old origin couplings keep their own
ledger; a current rare factor does not reset them.

\section{A checked scalar attachment for the full inherited interaction}

This is the most direct limited transfer result of this note.
It uses the actual full $V(X^-)$ in the selected R-II body, rather than
assigning a clean $O(g_e)$ norm to the entire inherited interaction.

Fix the positive component operation $T'_e(X)$ in R-II (1.71), including
its original internal history sum, before connected exponentiation.
Write the literal inherited exponent term as $W_z=V(X^-,z)$,
absorbing its original sign into $W$ if necessary. On the real reference
body,
\begin{equation}
 \dd\mu_{1,\omega}=e^{W_0}\dd\nu_\omega.
 \label{eq:realv}
\end{equation}
Here $\nu_\omega$ retains the original positive Wilson and Gaussian
factors, selectors, constraints, measures, reference denominators and
all other original scalars. No normalization is changed.
Here $z$ ranges over the original paired outer domain of this positive
operation. Additional source-only wall parameters may be carried
uniformly when their incoming envelope holds on this same body.
The argument does not extend the body comparison to an independently
changed preparation-wall kernel or to a signed localized term.

The original value bound (1.69) and its conversion into the positive
history cost give
\begin{equation}
 |W_z|\le C_VD_\omega,
 \label{eq:vvolume}
\end{equation}
where $D_\omega$ is the actual old-volume/geometric cost ledger, with
$D_\omega\le C_D\mathcal G_\omega$ for the original
$\mathcal G_\omega$ in \eqref{eq:livet}. All original old-scale counts
and collars remain in this comparison; the count $T_\omega$ alone is
not an arbitrary-volume bound.
The original outer comparison bounds the quadratic-form variation
$\sigma_Q$ and the full inherited variation $W_z-W_0$ separately.
These are source-free facts on the original native outer domain; they
do not admit the new mark~\cite{rii,W}.

\begin{proposition}[Uniform full-interaction interpolation]
For $0\le t\le1$, set
\begin{equation}
 \dd\mu_{t,\omega}=e^{tW_0}\dd\nu_\omega,\qquad
 \sigma_{t,z}=\sigma_{Q,z}+t(W_z-W_0).
 \label{eq:interpolation}
\end{equation}
The original positive numerical calculation, with all fixed extra costs
combined, supplies
\begin{equation}
 \sup_{0\le t\le1}\sum_\omega\int
 e^{\theta E_\omega+\tau D_\omega}\dd\mu_{t,\omega}
 \le B_X^{(c)},\qquad
 \sup_{t,\omega,z}|\sigma_{t,z}|\le\eps_*<\log2
 \label{eq:uniformv}
\end{equation}
on that original domain, for fixed positive surviving constants and a
sufficiently weak original ceiling.
\end{proposition}
\begin{proof}
$W_0$ is real and
$e^{tW_0}\le e^{|W_0|}\le e^{C_VD_\omega}$.
The original positive proof already uses this absolute majorant for
$V(X^-)$. Thus the same majorant works for every $t$; the extra
$e^{\tau D_\omega}$ merely increases its fixed original cost coefficient.
Retaining $\theta E_\omega$ and rerunning the birth/waiting/renewal/merger
recurrence is exactly the source-free reserve argument of S86 and
Section 6. It does not require $|W|=O(g_e)$.
For the complex comparison,
$|\sigma_t|\le\eps_Q+\eps_V$. In dimension four the source bound
has size $C g_eR_e^6p_0(g_e)^2q_0(g_e)$, which is small for the original
logarithmic choices. Both variations and the entire inherited $W$ have
been retained. Here $p_0(g),q_0(g)$ are the original cutoff functions,
while $p_0$ in $P_0$ denotes the declared logarithmic exponent.
\end{proof}

The proposition concerns the full resummed inherited scalar. An arbitrary
subfamily, a selected localized vertex, or its anchor derivatives is
covered only after its own absolute value and variation majorants are
verified. A bound on a full sum alone does not establish those stronger
bounds.

\subsection{One raw rare factor, with the complete history covariance}

Define
\begin{align}
 T_{t,z}F&=\sum_\omega\int e^{\sigma_{t,z}}F_{\omega,z}\dd\mu_{t,\omega},
 \qquad Z_{t,z}=T_{t,z}1,\nonumber\\
 \mathcal C_{t,z}(F,G)&=T_{t,z}(FG)
       -\frac{T_{t,z}F\,T_{t,z}G}{Z_{t,z}}.
 \label{eq:rawconn}
\end{align}
If $M_t=\sum_\omega\mu_{t,\omega}1>0$, then
$|Z_{t,z}|\ge(2-e^{\eps_*})M_t$. Cauchy--Schwarz on this one body gives
\begin{equation}
 |\mathcal C_{t,z}(F,G)|
 \le C_{\eps_*}\norm F_{L^2(\mu_t)}\norm G_{L^2(\mu_t)},
 \quad C_{\eps}=e^{\eps}+\frac{e^{2\eps}}{2-e^{\eps}}.
 \label{eq:connineq}
\end{equation}
This is the comparison of~\cite[\S3]{CE}. Its history sum is inside
the body; the between-history conditional-mean term in
\eqref{eq:totalcov} is therefore included. Any still external history
variables remain external and require their own subsequent estimate.

From \eqref{eq:vvolume} and \eqref{eq:uniformv},
\begin{equation}
 \norm{W_z}_{L^2(\mu_t)}
 \le \frac{2C_V}{\mathrm e\tau}\sqrt{B_X^{(c)}},
 \label{eq:v2}
\end{equation}
because $\sup_{d\ge0}d^2e^{-\tau d}=(2/(\mathrm e\tau))^2$.
There has been no new rare bank assigned to this interaction.

\begin{theorem}[Conditional full-$V$ scalar three-row attachment]
Assume the actual stored mark on the same complete selected support and
the whole original simultaneous domain satisfies
\eqref{eq:envelope} and $Q\le C/u$, with the unchanged clean radius.
Then for every fixed $0<\delta<c$,
\begin{equation}
 \boxed{
 |\mathcal C_{t,z}(F,W)|
 \le Cg_e A_j\rho^2
 \left[r^s+ur^b+u^2r^2\right] B_X^{(c-\delta)}
 }
 \label{eq:mainpartial}
\end{equation}
uniformly in $0\le t\le1$ and the verified outer domain in
\eqref{eq:uniformv}. This is a scalar raw attachment to the
literal full inherited interaction, conditional on the stated actual
incoming admission and envelope.
\end{theorem}
\begin{proof}
For the degree-$m$ envelope term write
$c_m=A_j\rho^2r^{\gamma_m}u^{m-2}$, where
$(\gamma_2,\gamma_3,\gamma_4)=(s,b,2)$.
Apply the pointwise envelope before the $L^2$ estimate. Its contribution
to $\norm F_{L^2(\mu_t)}$ is at most $Cc_mu^{-m}\sqrt{B_X^{(c)}}$.
Equations \eqref{eq:connineq}--\eqref{eq:v2} bound the raw contribution
by $Cc_mu^{-m}B_X^{(c)}$. To extract $g_ec_m$, the needed cost is
\begin{equation}
 g_e^{-1}u^{-m}=\alpha^m x_e^{(m+1)/2},\qquad 2\le m\le4.
 \label{eq:sharpcost}
\end{equation}
Its largest current power is $x_e^{5/2}$. Equation
\eqref{eq:rarepower} pays the whole finite set with one reduction
$c\mapsto c-\delta$. Summing the three bounds proves
\eqref{eq:mainpartial}. Both $L^2$ roots use the same majorant and
produce one $B$, not two independent masses. No saved age is spent.
\end{proof}

The pointwise or sufficiently reserve-weighted hypothesis is essential:
a bare $L^2(\mu_1)$ estimate does not imply uniformity under
\eqref{eq:interpolation}. The factor $g_e$ in
\eqref{eq:mainpartial} comes from rarity, not a pointwise clean-vertex
bound for $V$. The scalar maximum in \eqref{eq:rarepower}, with
$m=5/2$, makes its constant explicit.

A genuinely common outside scalar $a(z)$ multiplies
\eqref{eq:mainpartial} once. It is not canceled. History-dependent
reference factors stay inside their corresponding $\mu_{t,\omega}$.
For a readout held fixed during the exponent interpolation,
\begin{equation}
 Z_{t,z}\,\partial_t\left(\frac{T_{t,z}F}{Z_{t,z}}\right)
 =\mathcal C_{t,z}(F,W_z).
 \label{eq:normalizedderiv}
\end{equation}
Dividing \eqref{eq:mainpartial} by $Z_{t,z}$ leaves the genuine $B/Z$
problem. The local complex lower bound is proportional to the actual
mass $M_t$ and does not compare that mass to its rare upper majorant.

\section{The remaining main lemma, stated without hidden replacements}

The new results settle parts of the numerical carrying mechanism. They
leave the following precise geometric/analytic assertion unproved.
For the actual $\Gamma^R_{\omega,\mathbf s}$, one needs the full joined
target family, including its base values, to satisfy
\eqref{eq:necessary} with the original spare margins. That condition is
only necessary. The branch $P_{\rm old}q_{\rm old}\Gamma^R$ must exist
analytically on the entire simultaneous domain, and the actual pair
$(W,J_{\rm phys}(W))$ must lie in $\mathcal C^{\rm spare}_{j,e}$.
One then still needs the saved packet norm
\begin{equation}
 \mathfrak N_I\!\left[
 R_j\bigl(P_{\rm old}q_{\rm old}\Gamma^R,
     J_{\rm phys}(P_{\rm old}q_{\rm old}\Gamma^R)\bigr)\right]
 \le\sum_\gamma a_{I,\gamma;j,e}Q^{D_{I,\gamma}}.
 \label{eq:missingnorm}
\end{equation}
Here $\mathfrak N_I$ is the \emph{actual} finite current/source-jet and
weighted absolute anchor norm. Each $a_{I,\gamma;j,e}$ retains the old
mesh, radius factors, lowered coupling grade, saved age and source
coefficient. Formula \eqref{eq:missingnorm} is a target, not a theorem
proved by writing it down.

There are three concrete remaining checks.
\begin{enumerate}
\item \textbf{Joined targets and base traces.} Bound the complete
$q_{\rm old}\Gamma^R$ relative to the saved target domain, including
the bases in \eqref{eq:base}, all boundary identity and prefix rows,
and every simultaneous wall. Apply the old response theorem only
after that admission. The stationary/row defects
\eqref{eq:offbase} and \eqref{eq:rowdefect} must then be bounded on
their actual adjoint supports.
\item \textbf{The unassigned part of the aggregate.} Identify which
actual read rows are continuously live owned records eligible for
\eqref{eq:aggregatemoment}. For the remaining exterior and artificial
joined rows, prove the marginal bound \eqref{eq:pinnedtarget} or a
uniform deterministic full-record bound. An old scale label is not
its ancestry. The directional-current, nonlinear, gauge and reference
responses stay in that sum.
\item \textbf{The complete differentiated attachment.} Establish
\eqref{eq:missingnorm} and the corresponding localized distinguished
vertex bounds, with the full inherited effective score and current
chain rule. The scalar full-$V$ theorem does not supply the actual
MT36 kernel/anchor norm. Whole-kernel weak derivatives retain sharp
boundary partners, the Casimir action row and independent absolute
normalizer responses.
\end{enumerate}

No R-I range/selector statement is used here beyond
R-II's references to it. In particular the R-I preparation range
(1.64)--(1.69) at $n=N$, selector and cutoff clauses are not
used directly here. This is a limitation of the cited input, distinct
from the mathematical need to connect that native range to the saved
physical/canonical domain. A native range statement alone would not
settle the latter.

\subsection{Status table}

\begin{tabularx}{\textwidth}{@{}p{0.32\textwidth}X@{}}
\toprule
Result & Verified scope and remaining condition\\
\midrule
Target and pinned-stencil identities & Exact on the admitted old record;
reconstruction preserves all its target jets.\\
Changed-base KKT response & Exact on a common admitted chart;
the full defect is still unbounded on the actual walls.\\
Protected first tangents & $r^2g_e$ and $r^3g_e$ follow on the same protected
chart from its assumed genuine response.\\
Owned aggregate moments & Actual continuous native-to-current live chains
and proved native row envelopes; no artificial or exterior ownership.\\
Full-core raw $L^2$ carrying & One remaining rare coefficient pays the
known coupling loss after actual saved-envelope admission.\\
Full inherited scalar attachment & Actual full $V$ and original positive
body; \eqref{eq:mainpartial} remains conditional on incoming admission.\\
Full LM47/MT36, normalized transfer & Not proved. Current jets, absolute
anchors, selected localized vertex and global $B/Z$ remain.\\
\bottomrule
\end{tabularx}

No new contact prescription, native-fiber equality, pressure comparison,
or source replacement is used in this note. The S94
ultraviolet-uniform comparison of~\cite[pp.~692, 618--621]{W} and other global interfaces are not closed
by these local raw estimates. The Yang--Mills existence proof remains
incomplete.

\section*{Verification and source record}

Exact finite computations verify the polynomial square and
full-core cancellations, the higher-jet countermodel, KKT signs and
projection identities with rational matrices, the off-base correction
example, the Hadamard coefficient-norm separation, and the same-body
complex raw covariance inequality on a nontrivial finite body. These
checks validate formulas and expose omitted terms. They do not test the
actual Wilson determining geometry, prove admission, or establish a
continuum limit. The analytic reserve proofs are given in the text.

Primary sources are cited at the equations used. Source-free positive
estimates are imported only on their original domains and weak-trajectory
assumptions. The bounds here do not establish the remaining claims
of~\cite[pp.~1222--1226]{W}. A conditional row of the status table is not a
completed proof.

\begin{thebibliography}{9}
\bibitem{W} Wilson-source estimates, volumes 1--10 (\texttt{lattice\_theory/source\_estimates} in this collection), cited by source page.
Particularly \S28.9.3, (8.3)--(8.4);
S86; MT37--48; LM21--23 and LM42--52; VC3--6; AE8--AE9; PE12, on
pp.~137--138, 614, 692--696, 774, 781--784, 932--935, 969--970, 1001 and 1174--1180.
\bibitem{CE} \emph{Connected estimates for selected operations}, \texttt{lattice\_theory/block\_and\_transfer/conditional\_estimates/connected\_estimates\_for\_selected\_operations.tex} in this collection.
Raw connected comparison on the same complete selected body, with its
incoming-admission and normalization limitations.
\bibitem{UA} \emph{Uniform averaging rows and finite-jet response}, \texttt{lattice\_theory/block\_and\_transfer/conditional\_estimates/uniform\_averaging\_rows\_and\_finite\_jet\_response.tex} in this collection.
Complete uniform old-target averaging and directional-current filters,
nonlinear and boundary remainders, finite path jets, and the unresolved
mixed incoming-representative estimate.
\bibitem{rii} T. Balaban, \emph{Large field renormalization. II.
Localization, exponentiation, and bounds for the R operation},
Commun. Math. Phys. \textbf{122} (1989), 355--392.
\href{https://doi.org/10.1007/BF01238433}{doi:10.1007/BF01238433}.
Used: (1.50)--(1.65), (1.69)--(1.75), and (1.79)--(1.89).
\bibitem{iii} T. Balaban, \emph{Convergent renormalization expansions for
lattice gauge theories}, Commun. Math. Phys. \textbf{119} (1988), 243--285.
Used: (2.11)--(2.15), pp.256--257.
\bibitem{variational} T. Balaban, \emph{The variational problem and
background fields in renormalization group method for lattice gauge
theories}, Commun. Math. Phys. \textbf{102} (1985), 277--309.
\href{https://doi.org/10.1007/BF01229381}{doi:10.1007/BF01229381}.
Native regularity (2), (9)--(10), stationary equations and response
formula (190), Proposition 9. Used only on admitted native domains.
\bibitem{ri} T. Balaban, \emph{Large field renormalization. I.
The basic step of the R operation}, Commun. Math. Phys. \textbf{122}
(1989), 175--202.
\href{https://doi.org/10.1007/BF01257412}{doi:10.1007/BF01257412}.
Cited for identification only: its full range/selector
clauses are not used here.
\end{thebibliography}
\end{document}
